% answers_short.tex —— 1980 年·全国卷（理科） · 简短答案速查
% 由 tools/make_short_answers.py 自 answers.tex 的 \gksolhead/\gkanswer 提取，
% 并与 1980/1980_数学真题_理_答案要点.md 逐题比对后生成，勿手改（改源 MD 后重跑）。

\begin{gknotice}
\textbf{使用说明}：本节只给最终答案便于快速核对，完整推导见解析版。答案为编写者独立推导，\textbf{非官方标准答案}。
\end{gknotice}

\gkpart{题目}
\begin{itemize}[leftmargin=2.6em]
  \item[\textbf{1.}] (1) $xy(x^2-3y^2)(x^2+3y^2)$；(2) $xy(x-\sqrt3y)(x+\sqrt3y)(x^2+3y^2)$；(3) $xy(x-\sqrt3y)(x+\sqrt3y)(x+\sqrt3\mathrm iy)(x-\sqrt3\mathrm iy)$
  \item[\textbf{2.}] 三圆心所成三角形的三边长依次为 $3$、$4$、$5$，故它是直角三角形
  \item[\textbf{3.}] 见解析版
  \item[\textbf{4.}] 见解析版
  \item[\textbf{5.}] 见解析版
  \item[\textbf{6.}] (1) $M=1$，$m=-1$，$T=\dfrac{10\pi}{k}$（$k>0$；按通行版本 $f(x)=\sin\left(\frac{kx}{5}+\frac{\pi}{3}\right)$）；(2) $k_{\min}=32$
  \item[\textbf{7.}] $\angle B=\arcsin\dfrac{\sqrt5-1}{2}$（即 $\arctan\sqrt{\dfrac{\sqrt5-1}{2}}$，约 $38.17^\circ$）
  \item[\textbf{8.}] 见解析版
  \item[\textbf{9.}] 圆心运动到 $\left(\dfrac{1-\sqrt5}{4},\,0\right)$（即 $x_0=-\dfrac{\sqrt5-1}{4}$）时
\end{itemize}
